\clearpage\thispagestyle{empty}
\pdfbookmark[0]{Sources for Pallas (December 1810)}{pallas1810dec-sources}
\begin{center}\textbf{Sources for Pallas (December 1810)}\end{center}
\noindent This AI translation follows the report in \href{https://gdz.sub.uni-goettingen.de/id/PPN236007467}{\textit{Werke}, vol.~VI} (Göttingen, 1874), pp.~320--322. The corresponding passages appeared in item LXII of \href{https://archive.org/download/bub_gb_YCI4AAAAMAAJ/page/n4.jpg}{\textit{Monatliche Correspondenz}, vol.~XXII}, December 1810, pp.~\href{https://archive.org/download/bub_gb_YCI4AAAAMAAJ/page/n594.jpg}{591}, \href{https://archive.org/download/bub_gb_YCI4AAAAMAAJ/page/n595.jpg}{592}, \href{https://archive.org/download/bub_gb_YCI4AAAAMAAJ/page/n596.jpg}{593}, \href{https://archive.org/download/bub_gb_YCI4AAAAMAAJ/page/n597.jpg}{594} and \href{https://archive.org/download/bub_gb_YCI4AAAAMAAJ/page/n598.jpg}{595}. The University of Michigan copy was digitized by Google and is available from the Internet Archive.
\par\smallskip\noindent The journal introduces a letter from Gauss dated 2 December. The collected text begins with the quotation, includes his prediction for 15 December, and ends with his reply to the report of the observation. In the table on p.~594 the journal gives the year 1810 and the declination $21^\circ\;58''7$. The collected edition gives 1809 and $22^\circ\;58''7$. The orbital elements on p.~592 also differ in the minutes of the angle in the eccentricity row and the seconds of inclination. The translation follows the collected edition. Translation and typesetting produced by AI.
\clearpage
